%!TEX root main.tex
\section{Hadamard automata and Hadamard-finite series}
\label{sec:Hadamard}

In this section, we study a weighted model of computation strongly related to polynomial automata
and the class of series its recognises.
%
We begin in~\cref{sec:Hadamard:preliminaries} with some mathematical preliminaries.
%
In~\cref{sec:Hadamard-finite series} we define Hadamard-finite series, generalising the recognisable series,
and present their closure properties.
%
Hadamard automata are defined and studied in~\cref{sec:Hadamard automata},
where they are shown to recognise precisely the Hadamard-finite series.
%
This leads to the main result of the section:
Hadamard-finite series constitute an effective prevariety~(\cref{thm:Hadamard prevariety})
and thus the commutativity problem is decidable for this model\-
(\cref{thm:Hadamard-finite series - decidable commutativity}).
%
We conclude in~\cref{sec:multivariate polyrec} with an application of the commutativity problem
to a class of multivariate number sequences defined by polynomial recursions~(\cref{thm:polyrec consistency}).

\subsection{Preliminaries}
\label{sec:Hadamard:preliminaries}

\subsubsection{Difference algebra}
\label{sec:difference algebra}

Let $\A = \tuple{A; 0, 1, {+}, {*}}$ be an \emph{algebra} (over the field of rational numbers),
that is a vector space equipped with an associative commutative bilinear product
with identity $1$~\cite[Ch.~3, §1]{Bourbaki:Algebra:I:1974}.
%
An algebra \emph{endomorphism} is a linear function $F : A \to A$ \st
%
\begin{align}
    \label{eq:endomorphism}
    \tag{endomorphism rule}
    F(1) = 1
        \qquad\text{and}\qquad
        F(\alpha * \beta) = F \alpha * F \beta,
        \quad\text{for all } \alpha, \beta \in A.
\end{align}
%
A \emph{difference algebra}~\cite{RMCohn:DifferenceAlgebra:1965,Levin:DifferenceAlgebra:2008}
is an algebra $\tuple{A; 0, 1, {+}, {*}, (F_j)_{1\leq j\leq d}}$
together with finitely many endomorphisms $F_1, \dots, F_d : A \to A$.
%
Unlike what is usually assumed in the difference algebra literature,
we do not require the $F_i$'s to pairwise commute.
%
Difference algebras will provide the semantic background for the rest of the section.

As an example, consider the algebra of multivariate polynomials
$\tuple{\poly \Q X; 0, 1, {+}, {\cdot}}$,
where $X = \tuple{X_1, \dots, X_k}$ is a tuple of commuting variables.
%
%When the variable names do not matter,
%we sometimes write $\poly \Q k$ instead of $\poly \Q X$.
%
An endomorphism $F$ acts on a polynomial $\alpha \in \poly \Q X$
just by evaluation $F \alpha = \alpha(F(X_1), \dots, F(X_k))$,
and thus it is uniquely defined once we have fixed $F(X_1), \dots, F(X_k) \in \poly \Q X$.
%
% This presentation choice will make it easier to compare homomorphisms
% with shuffles and infiltrations in the upcoming sections.
%
If we fix distinguished endomorphisms $F_1, \dots, F_d$,
then the algebra of polynomials acquires the structure of a difference algebra
$\tuple{\poly \Q X; 0, 1, {+}, {\cdot}, (F_j)_{1 \leq j \leq d}}$.
%
Other examples of difference algebras will be presented in \cref{sec:Hadamard algebra,sec:shift algebra}.

\subsubsection{Hadamard difference algebra}
\label{sec:Hadamard algebra}

We define a product operation on series
which will be central for the development of this section.
%
The \emph{Hadamard product}~\cite{Fliess:1974} of two series $f, g \in \series \Q \Sigma$,
denoted by $f \hadamard g$,
is defined element-wise:
For every $w \in \Sigma^*$,
$\coefficient w (f \hadamard g) := \coefficient w f \cdot \coefficient w g$.
%
For instance, if we have two polynomials $2 \cdot aab - c$ and $3 \cdot aab + b$ (in noncommuting variables),
their Hadamard product is $6 \cdot aab$.
%
This is an associative and commutative operation, it distributes over sum,
and the identity element is the series $\one = \sum_{w\in \Sigma^*} 1 \cdot w$ mapping every word to $1$,
giving rise to the \emph{Hadamard algebra} $\tuple{\series \Q \Sigma; \zero, \one, {+}, {\hadamard}}$.
%
The Hadamard product for series is reminiscent of the notion of \emph{intersection} in language theory
and of \emph{fully synchronous composition} in concurrency theory.

In order to make it easier to compare the Hadamard product
to the shuffle and infiltration products studied later,
we find it convenient to provide the following alternative,
coinductive definition
(\cf~\cite[Definition 8.1]{BasoldHansenPinRutten:MSCS:2017}):
%
For any two series $f, g \in \series \Q \Sigma$,
their Hadamard product $f \hadamard g$
is the unique series \st
%
\begin{align}
    \tag{${\hadamard}$-$\e$}
    \label{eq:hadamard:base}
    \coefficient \e {(f \hadamard g)}
        &= f_\e \cdot g_\e, \\
    \tag{${\hadamard}$-$\deriveleft a$}
    \label{eq:hadamard:step}
        \deriveleft a (f \hadamard g)
        &= \deriveleft a f \hadamard \deriveleft a g,
        \ \forall a \in \Sigma.
\end{align}
%
%Thus it gives rise to the \emph{Hadamard algebra} of series
%$\series \Q \Sigma_{\hadamard} = \tuple{\series \Q \Sigma, {+}, {\hadamard}}$.
%
Equation \cref{eq:hadamard:step} states that left derivatives are endomorphisms \wrt~the Hadamard product.
%
We thus obtain the \emph{Hadamard difference algebra} of series
%\begin{align*}
    $\series \Q \Sigma_{\hadamard} := \tuple{\series \Q \Sigma; \zero, \one, {+}, {\hadamard}, (\deriveleft a)_{a \in \Sigma}}$,
%\end{align*}
which will provide the semantic background
for the forthcoming Hadamard automata.
%
Right derivatives are also Hadamard algebra endomorphisms,
which will be useful later.
% which will be used later in the proof of~\cref{lem:Hadamard closure properties}.
%
\begin{restatable}{proposition}{lemRightDerivationHadamardEndo}
    \label{lem:right derivation Hadamard endomorphism}
    Right derivatives $\deriveright a$ (with $a \in \Sigma$) are Hadamard algebra endomorphisms.
    I.e., they are linear, preserve the multiplicative identity $\deriveright a \one = \one$,
    and commute with Hadamard product:
    %
    \begin{align}
        \tag{${\hadamard}$-$\deriveright a$}
        \label{eq:derive hadamard right}
        \deriveright a (f \hadamard g) = \deriveright a f \hadamard \deriveright a g,
        \quad\forall a \in \Sigma.
    \end{align}
\end{restatable}
%
% \begin{proof}
%     Linearity follows directly from the definition of $\deriveright a$.
%     Equation~\cref{eq:derive hadamard right} could be proved coinductively,
%     however a direct argument suffices since the Hadamard product acts element-wise:
%     For every $w \in \Sigma^*$ and $a \in \Sigma$,
%     %
%     \begin{align*}
%         \coefficient w {\deriveright a (f \hadamard g)}
%         &= \coefficient {wa} {(f \hadamard g)}
%         = \coefficient {wa} f \cdot \coefficient {wa} g
%         = \coefficient w \deriveright a f \cdot \coefficient w \deriveright a g
%         = \coefficient w (\deriveright a f \hadamard \deriveright a g). \qedhere
%     \end{align*}
% \end{proof}

\subsection{Hadamard-finite series}
\label{sec:Hadamard-finite series}

In this section we rely on the structure of the Hadamard algebra
to define a finitely presented class of series.
%
For series $f_1, \dots, f_k \in \series \Q \Sigma$,
let $\poly \Q {f_1, \dots, f_k}_{\hadamard}$ be the smallest Hadamard algebra containing $f_1, \dots, f_k$.
%
Hadamard algebras of this form are called~\emph{finitely generated},
with \emph{generators} $f_1, \dots, f_k$.
%
It is a \emph{difference Hadamard algebra}
if it is additionally closed under the endomorphisms $\deriveleft a$, $a \in \Sigma$.
%
\begin{definition}[Hadamard-finite series]
    \label{def:Hadamard finite}
    A series is \emph{Hadamard finite} if it belongs to a
    finitely generated difference Hadamard algebra.
\end{definition}
%
\noindent
The following lemma will be our working definition of Hadamard-finite series.
%
\begin{restatable}{lemma}{lemHadamardFiniteWorkingDefinition}
    \label{lem:Hadamard finite - working definition}
    A series $f$ is Hadamard finite iff there are generators $f_1, \dots, f_k$ \st:
    %
    \begin{enumerate}
        \item $f \in \poly \Q {f_1, \dots, f_k}_{\hadamard}$.
        %\Ie, there exists a polynomial $p \in \poly \Q {y_1, \dots, y_k}$
        %\st~$f = p(f_1, \dots, f_k)$.

        \item $\deriveleft a f_i \in \poly \Q {f_1, \dots, f_k}_{\hadamard}$
        for every $a \in \Sigma$ and $1 \leq i \leq k$.
        %
        \Ie, there are polynomials $p^{(a)}_i \in \poly \Q {y_1, \dots, y_k}$ \st
        %
        \begin{align}
            \label{eq:Hadamard finite equations}
            \deriveleft a f_i = p^{(a)}_i(f_1, \dots, f_k),
            \quad\text{for all } a \in \Sigma, 1 \leq i \leq k.
        \end{align}
    \end{enumerate}
    %
    Moreover, we can assume without loss of generality that $f$ equals one of the generators $f = f_1$.
\end{restatable}
%
\noindent
We call~\cref{eq:Hadamard finite equations}~a \emph{system of Hadamard-finite equations}.
%
Note that every choice for the constant terms $\coefficient \e {f_i} \in \Q$
extends to a unique series solution of~\cref{eq:Hadamard finite equations}.
%
Consequently, a Hadamard-finite series can be finitely presented by
the constant terms and the polynomials $p_i^{(a)} \in \poly \Q {y_1, \dots, y_k}$.
%
%
\begin{proof}
    If $f$ is Hadamard finite, then the two conditions hold by definition.
    \Wlg~we can take $f = f_1$ to be one of the generators
    since we can just add $f$ to the generators,
    without changing the validity of the two conditions.

    For the other direction, assume that there are generators $f_1, \dots, f_k$ with $f = f_1$ satisfying the two conditions.
    %
    We need to show that the finitely generated algebra $A := \poly \Q {f_1, \dots, f_k}_{\hadamard}$ is closed under left derivatives.
    %
    To this end, consider a series $g \in A$.
    There is a polynomial $p \in \poly \Q {y_1, \dots, y_k}$ \st~$g = p(f_1, \dots, f_k)$.
    %
    Take now the left derivative of both sides.
    Since they are homomorphisms of $A$, we have
    $\deriveleft a g = \deriveleft a (p(f_1, \dots, f_k)) = p(\deriveleft a f_1, \dots, \deriveleft a f_k)$.
    %
    By the second condition, $\deriveleft a f_1, \dots, \deriveleft a f_k \in A$,
    and thus $\deriveleft a g \in A$ since $A$ is an algebra, as required.
\end{proof}

\begin{remark}
    Hadamard-finite equations have been considered under the name of
    \emph{polynomial recurrent relations} by Sénizergues\-
    \cite[Definition 5]{Senizergues:CSR:2007},
    where they are claimed, without proof, to coincide
    with those recognised by deterministic pushdown automata of level 3\-
    \cite[Corollary 3]{Senizergues:CSR:2007},
    and to have a decidable equality problem\-
    \cite[Theorem 5]{Senizergues:CSR:2007}.
    %
    In the context of arbitrary semirings,
    they also appear under the name of \emph{Hadamard-polynomial equations}~\cite[Sec.~6]{KostolanyiMisun:TCS:2018}.
\end{remark}

\begin{remark}[Cauchy product and Cauchy-finite series]
    The definition of Hadamard-finite series (as well as the shuffle and infiltration-finite series from the next sections) follows a natural pattern.
    %
    In fact, \emph{every} algebra of series $\tuple{\series \Q \Sigma; 0, 1, {+}, {\product}}$
    gives rise to a corresponding class of $*$-finite series:
    A series $f \in \series \Q \Sigma$ is \emph{$*$-finite}
    if it belongs to a finitely generated sub-algebra closed under $\deriveleft a$ ($a \in \Sigma$).
    %
    For instance, taking ``$\product$'' to be the \emph{Cauchy product} (also known as \emph{concatenation product})
    we obtain the \emph{Cauchy-finite series},
    which have been studied in~\cite{Wechler:MFCS:1979} and shown to coincide with the class of algebraic series.
    %
    Note that Cauchy product is not commutative, unlike the products that we consider in this paper,
    and thus does not fall under the scope of the techniques that we develop.
\end{remark}

\begin{remark}
    When all the polynomials involved in the definition of a Hadamard-finite series have degree $\leq 1$
    (\ie, they are linear), we recover the notion of recognisable series from~\cref{sec:recognisable}.
\end{remark}

We recall some basic closure properties of Hadamard-finite series.
All of them follow straightforwardly from the definitions.

\begin{restatable}[Closure properties of Hadamard-finite series]{lemma}{lemHadamardClosureProperties}
    \label{lem:Hadamard closure properties}
    The class of Hadamard-finite series is an effective difference algebra,
    \ie, it contains $\zero$, $\one$, and it is effectively closed under scalar product,
    sum, Hadamard product, and left derivatives.
    %
    Moreover, Hadamard-finite series are effectively closed under right derivatives and left anti-derivatives.
    %
    % Additionally,
    % \begin{itemize}
    %     \item if $f$ is a Hadamard finite,
    %     then $\deriveright a f$ is effectively Hadamard finite, and
    %     %
    %     \item If $g \in \series \Q \Sigma$ is an anti-derivative
    %     of a tuple of Hadamard-finite series $f = \tuple{f_a}_{a \in \Sigma} \in \series \Q \Sigma^\Sigma$,
    %     then $g$ is effectively Hadamard finite.
    % \end{itemize}
\end{restatable}
%
\noindent
When we say that Hadamard-finite series are effectively closed under right derivatives
we mean that there is an algorithm that, given in input a finite presentation for $f$ and an input symbol $a \in \Sigma$,
produces in output a finite presentation for $\deriveright a f$.
%
A similar remark applies to the effectiveness of the other closure properties.
%
Closure under anti-derivatives will be used in conjunction to~\cref{lem:equality reduces to commutativity}
to show that the commutativity problem is at least as hard as the equality problem.
%
\begin{proof}
    The series $\zero$ and $\one$ are obviously Hadamard finite.
    %
    Effective closure under scalar product $c \cdot f$ (with $c \in \Q$), addition $f + g$,
    left derivatives $\deriveleft a f$, and left anti-derivatives is done as in~\cref{lem:recognisable closure properties}.
    %
    Regarding closure under Hadamard product $f \hadamard g$,
    it is easy to show that if $f_1, \dots, f_k$ and $g_1, \dots, g_\ell$ are generators for $f$, resp., $g$,
    then $f_1, \dots, f_k, g_1, \dots, g_\ell$ can be taken to generate $f \hadamard g$.
    %
    % Since $\deriveleft a \zero = \zero$ and $\deriveleft a \one = \one$,
    % the series $\zero$ and $\one$ are both Hadamard finite.
    %
    We consider now closure under right derivatives.
    %
    Consider a Hadamard-finite series $f \in A$ for the finitely generated Hadamard algebra
    %
    \begin{align*}
        A := \poly \Q {f_1, \dots, f_k}_{\hadamard}
    \end{align*}
    %
    closed under left derivatives.
    %
    Fix an input symbol $a \in \Sigma$.
    %
    We consider new generators $\deriveright a f_i$, for every $1 \leq i \leq k$,
    yielding the finitely generated Hadamard algebra 
    %
    \begin{align*}
        B := \poly \Q {\deriveright a f_1, \dots, \deriveright a f_k}_{\hadamard}.
    \end{align*}
    %
    By~\cref{lem:Hadamard finite - working definition},
    the proof is concluded by the following two claims.
    %
    \begin{claim*}
        $\deriveright a f \in A$.
    \end{claim*}
    %
    \begin{claimproof}
        Since $f$ is in the Hadamard algebra generated by the $f_i$'s
        and right derivative is an endomorphism by~\cref{lem:right derivation Hadamard endomorphism},
        $\deriveright a f$ is in the Hadamard algebra generated by the $\deriveright a f_i$'s,
        \ie, $\deriveright a f \in A$.
    \end{claimproof}
    %
    \begin{claim*}
        For every $b \in \Sigma$ and $1 \leq i \leq k$,
        $\deriveleft b \deriveright a f_i \in A$.
    \end{claim*}
    %
    \begin{claimproof}
        Since left and right derivatives commute by~\cref{lem:derive left right commutativity},
        we have $\deriveleft b \deriveright a f_i = \deriveright a \deriveleft b f_i$.
        %
        But $\deriveleft b f_i$ is in the Hadamard algebra generated by the $f_i$'s
        and since right derivative is an endomorphism by~\cref{lem:right derivation Hadamard endomorphism},
        $\deriveright a \deriveleft b f_i \in A$.
    \end{claimproof}
\end{proof}

In this section we have defined the class of Hadamard-finite series in a semantic fashion,
based on which we have shown effective closure properties in~\cref{lem:Hadamard closure properties}.
%
In order to show that Hadamard-finite series constitute an effective prevariety,
it remains to show that the zeroness problem is decidable for them.
%
In order to achieve this, we introduce in the next section a syntactic presentation of the same class
via a corresponding notion of automata.

\subsection{Hadamard automata}
\label{sec:Hadamard automata}

We define the central computation model of the section.

\paragraph{Syntax.}

A \emph{Hadamard automaton} is a tuple $\AA = \tuple{\Sigma, X, F, \Delta}$
where $\Sigma$ is a finite \emph{input alphabet},
$X = \set{X_1,  \dots X_k}$ is a finite set of \emph{nonterminal symbols},
$F : X \to \Q$ is the \emph{output function} (or \emph{initial condition}),
and $\Delta : \Sigma \to X \to \poly \Q X$ is the \emph{transition function}.
%
A Hadamard automaton is \emph{linear} if $\Delta_a X_i$ is a polynomial of degree $\leq 1$
for all $a \in \Sigma$ and $1 \leq i \leq k$.

\paragraph{Semantics.}

A \emph{configuration} of a Hadamard automaton is a polynomial $\alpha \in \poly \Q X$.
%
For every input symbol $a \in \Sigma$,
the transition function $\Delta_a : X \to \poly \Q X$
is extended to a unique endomorphism $\Delta_a : \poly \Q X \to \poly \Q X$ of the polynomial algebra of configurations.
In other words, $\Delta_a$ acts by evaluation:
$\Delta_a(\alpha) = \alpha(\Delta_a(X_1), \dots, \Delta_a(X_k))$.
%
Finally, transitions from single letters
are extended to all finite input words homomorphically:
For every configuration $\alpha \in \poly \Q X$,
input word $w \in \Sigma^*$, and letter $a \in \Sigma$,
we have $\Delta_\e \alpha := \alpha$
and $\Delta_{a \cdot w} \alpha := \Delta_w \Delta_a \alpha$.
%
We have all ingredients to define the \emph{semantics} of a Hadamard automaton $\AA$,
which is a function $\Asem \AA \_ : \poly \Q X \to \series \Q \Sigma$
mapping a configuration $\alpha \in \poly \Q X$ to the series $\Asem \AA \alpha \in \series \Q \Sigma$ \st
%
\begin{align}
    \label{eq:semantics}
    \coefficient w {(\Asem \AA \alpha)} &:= F \Delta_w \alpha,
        \quad \text{for all } w \in \Sigma^*.
\end{align}
%
In other words, from the initial configuration $\alpha$ the automaton reads $w$ and goes to configuration $\Delta_w \alpha$,
from which the output function $F$ is applied to produce the output value $F \Delta_w \alpha \in \Q$.
% 
For this last step, $F$ is extended from nonterminals to configurations homomorphically
$F \beta = \beta(F X_1, \dots, F X_k)$. % (i.e., polynomial evaluation).
%
The series \emph{recognised} by the automaton $\AA$ is $\Asem \AA {X_1}$.

\begin{remark}
    The class of series recognised by linear Hadamard automata
    coincide with Schützenberger's recognisable series~\cite{Schutzenberger:IC:1961},
    which have been discussed in~\cref{sec:recognisable}.
\end{remark}

\begin{example}
    \label{ex:Hadamard automaton 1}
    Revisiting~\cref{ex:intro:1} in the syntax of Hadamard automata,
    we have a single nonterminal $X = \set{A}$,
    output function $F_c$ defined by $F_c A := c$ ($c \in \Q$),
    and transitions by
    %
    \begin{align*}
        \Delta_{a_1} A := A^2
            \quad\text{and}\quad
                \Delta_{a_2} A := 1 - A^2.
    \end{align*}
    %
    For instance, when $c = 0$ starting from the initial configuration $A$ and reading input $u = a_1 a_2$
    we obtain configuration $\Delta_{a_2} \Delta_{a_1} A = \Delta_{a_2} A^2 = (1 - A^2)^2$
    and thus $\Asem \AA A_u = 1$.
    %
    Reading $v = a_2 a_1$ yields a different configuration
    $\Delta_{a_1} \Delta_{a_2} A = 1 - A^4$,
    however $\Asem \AA A_v = 1$.
    %
    In fact, $\Asem \AA A_w$ is the count modulo $2$ of the number of $a_2$'s in $w$,
    and thus $\Asem \AA A$ is a commutative series.

    It can be checked that also the output functions $F_1, F_{-1}$
    defined by $F_1 A := 1$, resp., $F_{-1} A := -1$
    give rise to commutative series.
    %
    However, no other choice does.
    %
    For instance, $F_2$ defined by $F_2 A := 2$ does not yield a commutative series,
    since $\Asem \AA A_{a_1 a_2} = (1 - 2^2)^2 = 9$,
    but $\Asem \AA A_{a_2 a_1} = 1 - 2^4 = -15$.
\end{example}

A Hadamard automaton endows the algebra of configurations
with the structure of a difference algebra of polynomials
$\tuple{\poly \Q X; 0, 1, +, \cdot, (\Delta_a)_{a \in \Sigma}}$.
%
This allows us to connect the difference algebra of configurations
with the difference algebra of Hadamard series they recognise.
% In the next lemma we show that the semantics function connects
% the difference algebra of configurations with the difference algebra of series.
%
\begin{restatable}[Properties of the semantics]{lemma}{lemPropertiesofSemantics}
    \label{lem:Hadamard automata - properties of the semantics}
    The semantics of a Hadamard automaton is a \emph{homomorphism} from
    the difference algebra of polynomials
        $\tuple{\poly \Q X; 0, 1, +, \cdot, (\Delta_a)_{a \in \Sigma}}$
    to the difference Hadamard algebra of series 
        $\tuple{\series \Q \Sigma; \zero, \one, +, \hadamard, (\deriveleft a)_{a \in \Sigma}}$.
    %
    In other words, $\Asem \AA 0 = \zero$, $\Asem \AA 1 = \one$,
    and, for all $\alpha, \beta \in \poly \Q X$,
    %
    \begin{align}
        \label{eq:hadamard:sem:scalar-product}
        \Asem \AA {c \cdot \alpha}
            &= c \cdot \Asem \AA \alpha 
            && \forall c \in \Q, \\
        \label{eq:hadamard:sem:sum}
        \Asem \AA {\alpha + \beta}
            &= \Asem \AA \alpha + \Asem \AA \beta, \\
        \label{eq:hadamard:sem:product}
            \Asem \AA {\alpha \cdot \beta}
            &= \Asem \AA \alpha \hadamard \Asem \AA \beta, \\
        \label{eq:hadamard:sem:derivation}
        \Asem \AA {\Delta_a \alpha}
            &= \deriveleft a (\Asem \AA \alpha)
            && \forall a \in \Sigma.
    \end{align}
\end{restatable}
\begin{proof}
    Property $\Asem \AA 0 = \zero$ is trivial,
    and $\Asem \AA 1 = \one$ holds by definition:
    $\coefficient w \Asem \AA 1 = F \Delta_w 1 = F 1 = 1$.
    %
    Property~\cref{eq:hadamard:sem:derivation}
    holds by definition of the semantics.
    %
    Properties~\cref{eq:hadamard:sem:scalar-product} and~\cref{eq:hadamard:sem:sum} follow from linearity of $F$ and $\Delta_w$.
    %
    Finally, for property~\cref{eq:hadamard:sem:product}, we argue coinductively
    (the use of the coinductive assumption is marked by ``!''):
    %
    \begin{align*}
        \coefficient \e( \Asem \AA {\alpha \cdot \beta})
        &= F (\alpha \cdot \beta)
        = F \alpha \cdot F \beta
        = \coefficient \e (\Asem \AA) \alpha \cdot \coefficient \e (\Asem \AA \beta)
        = \coefficient \e (\Asem \AA \alpha \hadamard \Asem \AA \beta), \text{ and } \\
        %
        \deriveleft a (\Asem \AA {\alpha \cdot \beta})
            &= \Asem \AA {\Delta_a (\alpha \cdot \beta)}
            = \Asem \AA {\Delta_a \alpha \cdot \Delta_a \beta}
            \stackrel ! = \Asem \AA {\Delta_a \alpha} \hadamard \Asem \AA {\Delta_a \beta}
            = \deriveleft a \Asem \AA \alpha \hadamard \deriveleft a \Asem \AA \beta
            = \deriveleft a (\Asem \AA \alpha \hadamard \Asem \AA \beta). \qedhere
    \end{align*}
\end{proof}

\begin{remark}[Hadamard vs.~polynomial automata]
    \label{rem:Hadamard and polynomial automata}
    Hadamard automata are closely related to~\emph{polynomial automata}~\cite{BenediktDuffSharadWorrell:PolyAut:LICS:2017},
    a common generalisation of weighted automata~\cite{Schutzenberger:IC:1961}
    and vector addition systems~\cite{Mayr:STOC:1981}.
    %
    A state in a polynomial automaton is a tuple of rational numbers,
    which upon reading an input letter is modified by applying a polynomial update function;
    at the end of the run a polynomial output function is applied to produce the final weight.
    %
    To compare, Hadamard automata start from an initial polynomial
    and modify it by applying polynomial endomorphisms;
    at the end of the computation the state polynomial is evaluated on a particular tuple of rational numbers.
    %
    The two models are equivalent, up to reversal of the input:
    A series $f$ is recognised by a Hadamard automaton
    iff its reversal $f^R$ is recognised by a polynomial automaton.
    %
    %L: reversal already defined
    %(The \emph{reversal} of $f$ is the series $f^R$ mapping $a_1 \cdots a_n$ to $f(a_n \cdots a_1)$.)
    %
    We show this equivalence in detail in~\cref{app:polynomial automata}.
    
    Zeroness for polynomial automata is decidable with Ackermann complexity
    \-\cite[Corollary 1]{BenediktDuffSharadWorrell:PolyAut:LICS:2017},
    by an algorithm based on Hilbert's finite basis theorem
    originating in the work of Novikov and Yakovenko~\cite{NovikovYakovenko:1999}.
    %
    Since zeroness is invariant under reversal,
    this result also applies to Hadamard automata.
    %
    Since commutativity is invariant under reversal,
    our results also apply to polynomial automata.

    In the special case of an input alphabet with just one letter $\Sigma = \set a$,
    Hadamard (and polynomial) automata specialise to \emph{polynomial recursive sequences}\-
    \cite{Cadilhac:Mazowiecki:Paperman:Pilipczuk:Senizergues:ToCS:2021}.
    %
    They have also been considered in~\cite{BorealeGorla:CONCUR:2021} in a coalgebraic setting,
    where zeroness is shown to be decidable by the same algorithm as for polynomial automata above
    specialised to $\Sigma = \set a$~\cite[Theorem 4.1]{BorealeGorla:CONCUR:2021},
    however without complexity analysis.
    %
    It is currently an open problem to find an algorithm deciding the equality problem for polynomial recursive sequences
    of elementary complexity~\cite{ClementeDontenBuryMazowieckiPilipczuk:STACS:23}.

    Over more general commutative semirings,
    % (which are more general than the field of rational numbers considered in this paper),
    polynomial automata have been introduced as
    \emph{alternating weighted automata}~\cite{KostolanyiMisun:TCS:2018},
    further studied in~\cite{Grabolle:EPTCS:2021}
    (albeit without algorithmic results).
\end{remark}

The following lemma states that Hadamard automata recognise precisely the Hadamard-finite series.
%
Its proof relies on~\cref{lem:Hadamard automata - properties of the semantics,lem:Hadamard finite - working definition}.
%
\begin{restatable}{lemma}{lemHadamardAutomataAndSeries}
    \label{lem:Hadamard automata and series}
    A series is recognised by a Hadamard automaton if, and only if,
    it is Hadamard finite.
\end{restatable}
\begin{proof}
    % The proof is a direct consequence of the definitions.
    %
    For the ``only if'' direction, let $f = \Asem \AA {X_1}$
    be recognised by a Hadamard automaton $\AA =\tuple{\Sigma, X, F, \Delta}$
    with nonterminals $X = \tuple{X_1, \dots, X_k}$.
    %
    We show that $f$ is Hadamard finite
    by applying~\cref{lem:Hadamard finite - working definition}.
    %
    Consider the Hadamard algebra $A := \poly \Q {f_1,\dots, f_k}_{\hadamard}$
    finitely generated by series $f_i := \Asem \AA {X_i}$ for all $1 \leq i \leq k$.
    %
    Clearly $f \in A$.
    %
    Now consider generator $f_i$ and input symbol $a \in \Sigma$,
    and we need to show $\deriveleft a f_i \in A$.
    %
    By~\cref{lem:Hadamard automata - properties of the semantics},
    the semantics is a homomorphism of difference algebras, and thus
    %
    $\deriveleft a f_i
        = \deriveleft a {(\Asem \AA {X_i})}
        = \Asem \AA {\Delta_a X_i}
        = (\Delta_a X_i)(\Asem \AA {X_1}, \dots \Asem \AA {X_k})
        = (\Delta_a X_i)(f_1, \dots, f_k) \in A$,
    as required.

    For the ``if'' direction, let $f \in \series \Q \Sigma$ be Hadamard finite.
    %
    By~\cref{lem:Hadamard finite - working definition},
    there are series $f_1, \dots, f_k \in \series \Q \Sigma$ with $f = f_1$
    generating a Hadamard algebra $A = \poly \Q {f_1, \dots, f_k}_{\hadamard}$
    \st, for every generator $f_i$ and input symbol $a \in \Sigma$,
    there is a polynomial $p_i^a \in \poly \Q {X_1, \dots, X_k}$
    \st~$\deriveleft a f_i = p_i^a(f_1, \dots, f_k)$.
    %
    We have all ingredients to build a Hadamard automaton recognising $f_1$.
    %
    Consider the automaton $\AA = \tuple{\Sigma, X, F, \Delta}$
    with nonterminals $X := \tuple{X_1, \dots, X_k}$,
    output mapping $F(X_i) := \coefficient \e f_i$
    and transitions
    %
    \begin{align*}
        \Delta_a(X_i) := p_i^a(X_1, \dots, X_k),
        \quad \text{for all } 1 \leq i \leq k \text{ and } a \in \Sigma.
    \end{align*}
    %
    The proof is concluded by showing $\Asem \AA {X_i} = f_i$ for every $1 \leq i \leq k$.
    We argue coinductively.
    First, the two series have the same constant term by construction,
    %
    $\coefficient \e {(\Asem \AA{X_i})} = F \Delta_\e X_i = F X_i = \coefficient \e f_i$.
    %
    Second, the tails agree
    %
    \begin{align*}
        \deriveleft a {(\Asem \AA {X_i})}
        &= \Asem \AA {\Delta_a X_i} = 
            &&\text{(by~\cref{lem:Hadamard automata - properties of the semantics})} \\
        &= \Asem \AA {p_i^a(X_1, \dots, X_k)} = 
            &&\text{(def.~of $\Delta_a$)} \\
        &= p_i^a(\Asem \AA {X_1}, \dots, \Asem \AA {X_k}) = 
            &&\text{(by~\cref{lem:Hadamard automata - properties of the semantics})} \\
        &= p_i^a(f_1, \dots, f_k) = 
            &&\text{(by~coinduction)} \\
        &= \deriveleft a f_i.
            &&\text{(def.~of $p_i^a$)} \qquad \qedhere
    \end{align*}
\end{proof}

\begin{example}
    \label{ex:Hadamard finite equations 1}
    We illustrate~\cref{lem:Hadamard automata and series}
    on the Hadamard automaton from~\cref{ex:Hadamard automaton 1}.
    Let $f := \Asem \AA A$.
    %
    The corresponding Hadamard-finite equations are
    %
%    \begin{align*}
    $\deriveleft {a_1} f = f^2$ and $\deriveleft {a_2} f = \one - f^2$.
    % \end{align*}
    %
    (The square operation is to be understood as Hadamard square,
    \ie, $f^2 = f \hadamard f$;
    recall that $\one = \sum_{w \in \Sigma^*} 1 \cdot w$ is the Hadamard identity.)
    %
    It follows that $\poly \Q {f}_{\hadamard}$
    is a difference Hadamard algebra, and thus $f$ is Hadamard finite.
\end{example}

\begin{theorem}
    \label{thm:Hadamard prevariety}
    The class of Hadamard-finite series is an effective prevariety.
\end{theorem}
\begin{proof}
    The two closure conditions~\cref{prevariety:A} and~\cref{prevariety:B}
    follow from~\cref{lem:Hadamard closure properties}.
    %
    By~\cref{lem:Hadamard automata and series} and~\cref{rem:Hadamard and polynomial automata},
    the equality problem for Hadamard-finite series
    is (effectively) equivalent to the equality problem for polynomial automata,
    and the latter is decidable by~\cite[Corollary 1]{BenediktDuffSharadWorrell:PolyAut:LICS:2017}.
\end{proof}

We now have all ingredients to prove our main result for Hadamard-finite series.
%
\begin{theorem}
    \label{thm:Hadamard-finite series - decidable commutativity}
    The commutativity problem for Hadamard-finite series (equivalently, Hadamard and polynomial automata) is Ackermann-complete.
\end{theorem}
\begin{proof}
    Thanks to~\cref{thm:commutativity for effective prevarieties} and \cref{thm:Hadamard prevariety},
    the commutativity problem for Hadamard-finite series is decidable.
    %
    In fact, we have reduced the commutativity problem over alphabet $\Sigma$
    to $\card \Sigma^2 + \card \Sigma$ equivalence queries (\cf~\cref{lem:commutativity}).
    %
    On the other hand, by~\cref{lem:Hadamard closure properties},
    Hadamard-finite series are effectively closed under anti-derivatives,
    and thus by~\cref{lem:equality reduces to commutativity}
    the equivalence problem for Hadamard-finite series
    reduces to the commutativity problem.
    %
    Moreover, this reduction can be done efficiently, \ie, in polynomial time.
    %
    Since the complexity of equivalence for polynomial automata is Ackermann-complete\-
    \cite[Theorem 1 + Corollary 1]{BenediktDuffSharadWorrell:PolyAut:LICS:2017},
    the same is true for commutativity of Hadamard-finite series,
    Hadamard and polynomial automata.
\end{proof}

\subsection{Application: Multivariate polynomial recursive sequences}
\label{sec:multivariate polyrec}

We introduce a class of multivariate recursive sequences
generalising the univariate \emph{polynomial recursive sequences} (in short, \emph{polyrec})
from~\cite{Cadilhac:Mazowiecki:Paperman:Pilipczuk:Senizergues:ToCS:2021}.
%
We will make a crucial use of~\cref{thm:Hadamard-finite series - decidable commutativity}
to argue that this class can be presented \emph{effectively} by polynomial recursive equations.
%
In other words, we will show that polynomial recursive equations are a \emph{decidable syntax} for polyrec sequences:
%
There is an algorithm which, given in input a system of polynomial recursive equations and an initial condition,
decides whether they actually define a (polyrec) sequence at all.
%
The lack of an effective presentation
has been an obstacle to the study of a multivariate generalisation of polyrec sequences.

In~\cref{sec:shift algebra} we begin with some preliminaries about multivariate sequences,
in~\cref{sec:polyrec} we define multivariate polyrec sequences,
whose basic properties we study in ~\cref{sec:polyrec properties}.
%
Finally in~\cref{sec:consistency of polyrec equations} we use~\cref{thm:Hadamard-finite series - decidable commutativity}
to argue that their syntax is effective.

\subsubsection{Difference sequence algebra}
\label{sec:shift algebra}

Fix a dimension $d \in \N_{\geq 1}$.
%
For every $1 \leq j \leq d$,
by $e_j$ we denote the $j$-th unit vector in $\N^d$.
%
A \emph{multivariate sequence} is a function $f : \N^d \to \Q$.
%
We equip the set of sequences with the structure of a difference algebra
$\tuple{\N^d \to \Q; \zero, \one, {+}, {\cdot}, (\shift j)_{1 \leq j \leq d}}$,
where $\zero$, $\one$, scalar product $c \cdot f$ ($c \in \Q$),
sum $f + g$, and Hadamard product $f \cdot g$ are defined element-wise.
%
Note that we overload the notation for the scalar and Hadamard product.
%
Moreover, the \emph{$j$-th left shift} endomorphism $\shift j$ maps a sequence $f : \N^d \to \Q$
to the sequence $\shift j f$ \st
%
\begin{align*}
    (\shift j f)(n) := f(n + e_j),
    \quad \text{for all } n \in \N^d.
\end{align*}
%
Unlike the difference algebra of Hadamard series from~\cref{sec:Hadamard algebra},
the endomorphisms $\shift j$'s pairwise commute
$\shift j \circ \shift h = \shift h \circ \shift j$.
%which we call the \emph{shift algebra}.

\subsubsection{Polyrec sequences}
\label{sec:polyrec}

A $d$-variate sequence $f : \N^d \to \Q$ is \emph{polyrec}
if there exist $k \in \N_{\geq 1}$,
auxiliary sequences $f_1, \dots, f_k : \N^d \to \Q$ with $f = f_1$,
and polynomials $p^{(j)}_i \in \poly \Q {y_1, \dots, y_k}$ for all $1 \leq i \leq k$ and $1 \leq j \leq d$, \st
%
\begin{align}
    \label{eq:polyrec}
    \shift j f_i
        = p^{(j)}_i(f_1, \dots, f_k),
        \quad \forall 1 \leq i \leq k, 1 \leq j \leq d.
\end{align}
%
%for all $1 \leq i \leq k$ and $1 \leq j \leq d$.
%
In other words, for each $f_i$ and coordinate $1 \leq j \leq d$
there is a distinct recursive equation
specifying how future values can be computed from the current one.
%
We call systems of the form~\cref{eq:polyrec} \emph{systems of polyrec equations}.
%
Notice that the definition of a polyrec sequence contains the \emph{promise} that $f_1, \dots, f_k$ exist solving the polyrec equations,
however it is not trivial to establish whether solutions of polyrec equations exist,
and thus it is not clear whether polyrec sequences have an effective syntax.
%
We will return to this important issue in~\cref{sec:consistency of polyrec equations}.

The univariate case $d = 1$ corresponds to the polyrec sequences
from~\cite{Cadilhac:Mazowiecki:Paperman:Pilipczuk:Senizergues:ToCS:2021},
a rich class containing the linear recursive sequences (such as the Fibonacci numbers),
as well as fast growing sequences such as $2^{2^n}$.
%
Notice that in this case every initial condition extends (uniquely) to a solution of a system of polyrec equations~\cref{eq:polyrec},
and thus polyrec equations are an effective syntax for univariate polyrec sequences.
%
%L: this was said before
%The zeroness and equivalence problems for univariate polyrec sequences are decidable
%(with Ackermann complexity~\cite{BenediktDuffSharadWorrell:PolyAut:LICS:2017}),
%and whether there is an elementary algorithm is currently an open problem
%(\cf~\cite{ClementeDontenBuryMazowieckiPilipczuk:STACS:23}).
%
In order to develop intuition, we show below some examples of polyrec sequences.
%
\begin{example}
    We begin with a simple example, univariate and linear.
    %
    The \emph{Fibonacci sequence} $F : \N \to \Q$
    is well-known to satisfy the linear recursion $F(n+2) = F(n+1) + F(n)$.
    %
    We can turn this into the polyrec format by introducing an auxiliary sequence $G : \N \to \Q$ \st\-
    %
    $\shift 1 F = F + G$ and $\shift 1 G = F$.
    %
    % In fact, the difference equations above are even linear (with constant coefficients),
    % and thus the Fibonacci sequence is a linear recursive sequence.
    In this way, one shows that all linear recursive sequences are polyrec.
\end{example}

\begin{example}
    \label{ex:polyrec 1}
    We now consider a bivariate example $f : \N^2 \to \Q$.
    %
    For every $n_1, n_2 \in \N$, let 
    $f(n_1, n_2) := 2^{2^{n_1}} \cdot 2^{3^{n_2}}$.
    %
    Introducing auxiliary sequences $g(n_1, n_2) := 2^{2^{n_1}}$ and $h(n_1, n_2) := 2^{3^{n_2}}$
    we have polyrec equations
    %
    \begin{align*}
        \arraycolsep=1pt
        \begin{array}{rl}
            \shift 1 f &= f \cdot g, \\
            \shift 2 f &= f \cdot h^2,
        \end{array}
        \quad
        \begin{array}{rl}
            \shift 1 g &= g^2, \\
            \shift 2 g &= g,
        \end{array}
        \quad
        \begin{array}{rl}
            \shift 1 h &= h, \\
            \shift 2 h &= h^3.
        \end{array}
    \end{align*}
    % and thus $f, g, h$ are polyrec.
\end{example}

\begin{example}
    The bivariate factorial $f(n_1, n_2) := n_1! \cdot n_2!$ satisfies the recursion
    %
    \begin{align*}
        f(n_1+1, n_2) &= (\underline {n_1} + 1) \cdot f(n_1, n_2), \\
        f(n_1, n_2+1) &= (\underline {n_2} + 1) \cdot f(n_1, n_2).
    \end{align*}
    %
    This is not in the polyrec format~\cref{eq:polyrec},
    because of the two underlined occurrences of the index variables $n_1, n_2$.
    %
    Linear recursions with coefficients polynomial in $n_1, \dots, n_d$, as above,
    are called \emph{P-recursive}~\cite{Zeilberger:JMAA:1982}.
    %
    We can introduce auxiliary sequences $g(n_1, n_2) := n_1$ and $h(n_1, n_2) := n_2$
    and transform P-recursions into polyrec equations
    %
    \begin{align*}
        \arraycolsep=1pt
        \begin{array}{rl}
            \shift 1 f &= (g+\one) \cdot f, \\
            \shift 2 f &= (h+\one) \cdot f,
        \end{array}
        \quad
        \begin{array}{rl}
            \shift 1 g &= \one + g, \\
            \shift 2 g &= g,
        \end{array}
        \quad
        \begin{array}{rl}
            \shift 1 h &= h, \\
            \shift 2 h &= \one + h.
        \end{array}
    \end{align*}
    %
    A similar argument shows that $n_1! + n_2!$ is polyrec.
    %
    In fact, polyrec include all $P$-recursive definitions with a constant leading term
    (\cf~\cite[Proposition 3.6]{Lipshitz:D-finite:JA:1989}),
    also called \emph{monic $P$-recursive}\-
    \cite[page 5]{Buna-MargineanChevalShirmohammadiWorrell:POPL:2024}.
    %
    For instance, $n!$ is monic $P$-recursive since $(n+1)! = (n+1) \cdot n!$,
    but the natural recursion for the \emph{Catalan numbers} $C_n$, which is $\underline {(n+2)} \cdot C_{n+1} = (4n+2) \cdot C_n$, is not
    (the underlined term is non-constant).
    %
    In fact, the Catalan numbers are known not to be polyrec~\cite[Corollary 4.1]{Cadilhac:Mazowiecki:Paperman:Pilipczuk:Senizergues:ToCS:2021},
    and thus they are not monic P-recursive either.
\end{example}
%
\begin{remark}
    All the examples so far are sums of products of univariate polyrec sequences.
    %
    For examples not of this form,
    consider the polyrec sequences $(n_1 + n_2)!$ and $2^{2^{n_1 + n_2}}$.
    %
    This shows that the characterisation of multivariate \emph{linear} recursive sequences
    in~\cite[Proposition 2.1]{Karhumaki:TCS:1977} as sums of products of univariate linear recursive sequences
    does not generalise to polyrec.
\end{remark}
%
\noindent
More examples will be shown in~\cref{sec:consistency of polyrec equations}.

\subsubsection{Closure properties of polyrec sequences.}
\label{sec:polyrec properties}
%
Multivariate polyrec sequences are a well-behaved class of sequences.
%
They are closed under the algebra operations (scalar product, sum, and Hadamard product),
shifts $\shift j f$, sections, and diagonals.
%
The last two operations have not been defined yet.
%
A \emph{section} of a {$d$-variate} sequence $f : \N^d \to \Q$
is an {$e$-variate} sequence with $e < d$ which is obtained from $f$
by fixing one or more coordinates to a constant value~\cite[Definition 3.1]{Lipshitz:D-finite:JA:1989}.
%
For instance, if $f(n_1, n_2)$ is a bivariate sequence, then $g(n_1) := f(n_1, 7)$ is a section thereof,
obtained by fixing the second coordinate to be $n_2 := 7$.
%
Since sections are defined element-wise,
they commute with the operations of scalar product, sum, and Hadamard product,
and with shifts not involving the coordinates being fixed.
%
It follows that sections of polyrec sequences are polyrec.

\begin{example}
    For instance, if $\shift 1 f = \one - f^2$,
    then its section $g$ satisfies the initial condition $g(0) = f(0, 7)$
    and the same recursive equation $\shift 1 g = \one - g^2$,
    preserving the polyrec format
    (the information about $\shift 2 f$ is discarded when fixing the second coordinate).
\end{example}

A \emph{diagonal} of a {$d$-variate} sequence $f : \N^d \to \Q$
is an $e$-variate sequence with $e < d$ which is obtained from $f$
by requiring a subset of the coordinates to be equal
and then by projecting them to a single coordinate\-
\cite[Definition 2.6]{Lipshitz:D-finite:JA:1989}.
%
For instance, the sequence $h : \N^2 \to \Q$ \st\-
$h(n_1, n_3) := f(n_1, n_1, n_3)$, for all  $n_1, n_3 \in \N$,
is the diagonal of $f : \N^3 \to \Q$
obtained by identifying the first two coordinates.
%
Again, this is an element-wise operation,
and thus it commutes with the operations of scalar product, sum, Hadamard product,
and of shifts not involving the coordinates being identified.
%
Consequently, diagonals of polyrec sequences are polyrec.

\begin{example}
    For instance, if $\shift 1 f = \one - f^2$, $\shift 2 f = f^3$, and $\shift 3 f = \one + 2 \cdot f$
    then the diagonal $h$ satisfies $\shift 1 h = (\one - h^2)^3$ and $\shift 2 h = \one + 2 \cdot h$,
    preserving the polyrec format.
    %
    % Note that this is a conditional statement
    % and we do not claim that such an $f$ exists.
    %
    (By composing in the other order, we get $\shift 1 h = 1 - (h^3)^2$, but one equation suffices:
    If $f$ exists, then it will satisfy both equations.)
\end{example}

%
%Polyrec sequences constitute an effective difference algebra
%(this will follow from our isomorphism lemma~\cref{lem:Hadamard isomorphism lemma}).

\subsubsection{Consistency of polyrec equations}
\label{sec:consistency of polyrec equations}

So far we have defined a class of multivariate sequences, and we have argued that this class is robust
by presenting examples and closure properties.
%
However, we have sidestepped a crucial issue,
namely whether the equations~\cref{eq:polyrec} together with an initial condition
do actually define a sequence.
%
By the shape of the equations,
a given initial condition gives rise to \emph{at most one} solution.
%
However, whether a solution exists at all is nontrivial.
%
\begin{problem}[Polyrec consistency problem]
    \label{problem:polyrec consistency}
    The \emph{consistency problem} for polyrec equations
    takes as input a set of polyrec equations~\cref{eq:polyrec}
    together with an initial condition $c \in \Q^k$,
    and it amounts to decide whether there is a sequence solution $f \in (\N^d \to \Q)^k$ 
    extending the initial condition $f(0) = c$.
\end{problem}
%
\noindent
We emphasise that~\cref{problem:polyrec consistency} is about polyrec equations,
which may or may not represent a sequence,
and the question is finding out which one is the case.
%
We show that~\cref{problem:polyrec consistency} is decidable,
thus showing that the syntax of polyrec sequences~\cref{eq:polyrec} is \emph{effective}.
%
In the univariate case $d = 1$,
\emph{every} initial condition extends to a solution,
and thus the consistency problem is trivial.
%
However, in the multivariate case $d \geq 2$,
whether a system of equations~\cref{eq:polyrec}
actually defines a (tuple of) sequence(s)
involves an infinite number of constraints to be satisfied\-
(\cf~\cref{fig:polyrec constraints}).
%
In order to develop intuition, we provide some examples
illustrating the possible situations that may arise.

\begin{example}[All initial conditions extend to a solution]
    Let $d = 2$ and consider polyrec equations
    %
    \begin{align*}
        \shift 1 f = f^3
            \quad\text{and}\quad
                \shift 2 f = f^5.
    \end{align*}
    %
    The update functions $p(x) := x^3$ and $q(x) := x^5$
    commute identically $p(q(x)) = q(p(x)) = x^{15}$.
    %
    This is a very strong condition,
    implying that \emph{every} initial condition $f(0, 0) := c \in \Q$
    extends to a sequence solution,
    namely $f(n_1, n_2) = c^{3^{n_1} \cdot 5^{n_2}}$.
\end{example}

\begin{example}[Only some initial conditions extend to a solution]
    \label{ex:polyrec equations}
    Revisiting~\cref{ex:intro:2},
    let $d = 2$ and consider polyrec equations
    %
    \begin{align*}
        \shift 1 f = f^2
            \quad\text{and}\quad
                \shift 2 f = \one - f^2.
    \end{align*}
    %
    We claim that the set of initial values extending to a solution is $\set{-1, 0, 1}$.
    %
    Indeed, if we let $p(x) := x^2$ and $q(x) := 1 - x^2$,
    then we obtain the equation $p(q(x)) = q(p(x))$.
    %
    In other words, if the initial condition $x$ extends to a sequence solution,
    then it is necessarily a root of $p(q(x)) - q(p(x)) = (1-x^2)^2 - (1 - x^4) = 2 x^2(1-x)(1+x)$.
    %
    Over $\Q$, the latter polynomial has precisely three solutions $x \in \set{-1, 0, 1}$.
    % p(q(x)) = (1 -x^2)^2 = 1 - 2x^2 + x^4
    % q(p(x)) = 1 - x^4
    % => x = 0 or x^2 = 1
    %
    Moreover, in fact $x \in \set{-1, 0, 1}$ is also sufficient for the existence of a sequence solution,
    since this set is stable under $p(x)$ and $q(x)$:
    %
    \begin{center}
    \newcommand{\littlediagram}[4]{
        \tikzset{x=8mm, y=8mm, every node/.append style={font=\footnotesize}}
        \begin{tikzpicture}
            \draw[step=8mm] (0,0) grid (1,1);
            \foreach \i in {0,...,1}
                \foreach \j in {0,...,1}
                    \fill (\i,\j) circle (1pt);
            \draw[very thick, ->] (0,0) -- (1,0) node[midway, below] {$p$};
            \draw[very thick, ->] (0,0) -- (0,1) node[midway, left] {$q$};
            \draw[very thick, ->] (0,1) -- (1,1) node[midway, above] {$p$};
            \draw[very thick, ->] (1,0) -- (1,1) node[midway, right] {$q$};
            \fill[black]
                (0,0) circle (2pt) node[below] {$#1$}
                (1,0) circle (2pt) node[below] {$#2$}
                (0,1) circle (2pt) node[above] {$#3$}
                (1,1) circle (2pt) node[above] {$#4$};
            \end{tikzpicture}
        }
    \begin{tabular}{ccc}
        \littlediagram {-1} 1 0 0&
        \littlediagram 0 0 1 1&
        \littlediagram 1 1 0 0
    \end{tabular}
    \end{center}
    %
    Consequently, all and only $f(0, 0) \in \set{-1, 0, 1}$ extend to a solution.
\end{example}

% \begin{example}[Infinitely many initial conditions extend to a solution]
%     \label{ex:infinitely any initial conditions}
%     The equations from~\cref{ex:polyrec 1} admit a solution for every initial condition
%     $c := f(0, 0), d := g(0, 0), e := h(0, 0) \in \Q$ \st~$c = d \cdot e$.
%     In this case the unique solution is
%     $f(n_1, n_2) = d^{2^{n_1}} \cdot e^{3^{n_2}}$,
%     $g(n_1, n_2) = d^{2^{n_1}}$,
%     $h(n_1, n_2) = e^{3^{n_2}}$.
%     %
%     In this case we have infinitely many initial conditions extending to a solution.
%     \lorenzo{In fact all initial conditions extend to a solution. fix this}
% \end{example}

\begin{example}[No initial condition extends to a solution]
    Let $d = 2$ and consider the polyrec equations
    %
    \begin{align}
        \label{eq:polyrec example 2}
        \shift 1 f = f^2
            \quad\text{and}\quad
                \shift 2 f = f + \one.
    \end{align}
    %
    Let the initial condition be $f(0, 0) := c \in \Q$.
    From the first equation we have $f(1, 0) = c^2$,
    and from the second $f(0, 1) = c+1$.
    %
    But then we can obtain $f(1, 1)$ in two ways,
    %
    $c^2 + 1$ and $(c+1)^2 = c^2 + 2c + 1$,
    %
    implying $c = 0$ and thus $f(1, 1) = 1$.
    %
    But now we compute $f(2, 2)$ in two ways
    obtaining $1^2 + 1 = 2$ and $(1+1)^2 = 4$.
    This leads to a contradiction,
    and thus there is no initial condition extending to a solution of~\cref{eq:polyrec example 2}.
\end{example}

% perhaps not interesting
% \begin{example}[Exactly one initial condition extends to a solution]
%     Let $d = 2$ and consider the polyrec equations
%     %
%     \begin{align}
%         \label{eq:polyrec example 3}
%         \shift 1 f = f^2
%             \quad\text{and}\quad
%                 \shift 2 f = 2\cdot f.
%     \end{align}
%     %
%     Clearly the initial condition $f(0, 0) = 0$
%     extends to the trivial solution $f = 0$.
%     %
%     In fact, no other initial condition $c := f(0, 0) \neq 0$ extends to a solution:
%     Computing $f(1, 1)$ in two ways we obtain $2 \cdot c^2 = (2c)^2$,
%     which is impossible when $c \neq 0$.
% \end{example}

%We are thus left with the problem of deciding whether a given system of polyrec equations~\cref{eq:polyrec} admits a solution.
%When the initial condition $c \in \Q^k$ is given,
%this is the same as checking membership $c \in C$ in the algebraic variety $C$ from~\cref{rem:algebraic variety of initial conditions}.
%
%When the initial condition is not fixed, the task is to decide whether $C$ is nonempty.
\noindent
\cref{problem:polyrec consistency} is an instance of the more general problem
solvability in sequences for systems of polynomial difference equations,
not necessarily of the polyrec kind~\cref{eq:polyrec}.
%
The latter problem is undecidable in full generality, as we now recall.
%
\begin{remark}
    \label{rem:undecidability of polynomial difference equations}
    The more general problem of solvability in sequences
    for systems of polynomial difference equations
    is undecidable~\cite[Proposition 3.9]{PogudinScanlonWibmer:2020}.
    %
    Nonlinear bivariate systems of two unknowns $k = d = 2$ suffice.
    %
    The undecidable systems~\cite[Eq.~(5) on pg.~21]{PogudinScanlonWibmer:2020}
    in our notation can be written as follows:
    %
    \begin{align*}
        (f - 1)(f - 2) \cdots (f - N) &= 0, \\
        (g - 1)(g - 2) \cdots (g - N) &= 0, \\
        \prod_{i = 1}^\ell ((a_i - f)^2 + (b_i - \shift 2 f)^2 + (c_i - g)^2 + (d_i - \shift 1 g)^2) &= 0,
    \end{align*}
    %
    for suitable constants $N, a_1, \dots, a_\ell, b_1, \dots, b_\ell, c_1, \dots, c_\ell \in \N$,
    where for simplicity of notation, we have identified $n \in \N$ with $n \cdot \one$.
    %
    The equations above encode plane tiling problems:
    the first two constraints force $f, g : \N^2 \to \set {1, \dots, N}$ to range over a finite set,
    and the last equation encodes horizontal and vertical tiling constraints.
    % (we do not need the details of the encoding here).
    %
    None of the equations above is in the polyrec format~\cref{eq:polyrec}.
\end{remark}

Consequently, in order to solve~\cref{problem:polyrec consistency}
we need to exploit the special shape of polyrec equations.
%
To this end, we develop a connection between sequence solutions of polyrec equations~\cref{eq:polyrec}
and Hadamard-finite series.
%
We make the very simple observation
that number sequences and commutative series contain the same information.
%
For instance, the sequence $2^{n_1} 3^{n_2}$ is ``the same'' as the series that maps
$w \in \set{a_1, a_2}^*$ to $2^{\multiplicity {a_1} w} 3^{\multiplicity {a_2} w}$.
%
This is made precise by saying that the difference sequence algebra
$\tuple{\N^d \to \Q; \zero, \one, +, \cdot, (\shift j)_{1 \leq j \leq d}}$
and the difference Hadamard algebra of commutative series
$\tuple{\commseries \Q \Sigma; \zero, \one, +, \hadamard, (\deriveleft {a_j})_{1 \leq j \leq d}}$
(with $\Sigma = \set{a_1, \dots, a_d}$)
are isomorphic.
%
The isomorphism maps a commutative series $f \in \commseries \Q \Sigma$
to the number sequence $\widetilde f : \N^d \to \Q$
\st~$\widetilde f(n_1, \dots, n_d) := f(a_1^{n_1} \cdots a_d^{n_d})$ for every $n_1, \dots, n_d \in \N$.
%
This map preserves the difference algebra operations, which has pleasant consequences.

Given a system of polyrec equations~\cref{eq:polyrec},
we construct a system of Hadamard-finite equations
%
\begin{align}
    \label{eq:Hadamard-finite companion system}
    \deriveleft {a_j} g_i = p^{(j)}_i(g_1, \dots, g_k),
    \quad \forall 1 \leq i \leq k, 1 \leq j \leq d,
\end{align}
%
which we call the \emph{companion system} to~\cref{eq:polyrec}.
%
\begin{example}
    The polyrec equations from~\cref{ex:polyrec equations}
    $\set{\shift 1 f = f^2, \shift 2 f = \one-f^2}$
    yield the companion system from~\cref{ex:Hadamard finite equations 1}
    %
    $\set{\deriveleft {a_1} g = g^2, \deriveleft {a_2} g = \one - g^2}$.
    %
\end{example}
%
\noindent
Every initial condition extends to a \emph{unique} series solution $g = \tuple{g_1, \dots, g_k} \in \series \Q \Sigma^k$
of the companion system~\cref{eq:Hadamard-finite companion system},
which can be proved by induction on the length of words.
%
By definition, this solution is a tuple of Hadamard-finite series.
%
When all components $g_1, \dots, g_k$ are commutative,
we can extract corresponding number sequences
$f_1 := \widetilde {g_1}, \dots, f_k := \widetilde {g_k} : \N^d \to \Q$,
which are precisely the unique sequence solution
to the original polyrec equations~\cref{eq:polyrec}:
%
\begin{align*}
    \shift j {f_i}
   = \shift j \widetilde {g_i}
    = \widetilde {\deriveleft {a_j} g_i}
    = \widetilde {\left(p^{(j)}_i(g_1, \dots, g_k)\right)}
   = p^{(j)}_i(\widetilde {g_1}, \dots, \widetilde {g_k})
    = p^{(j)}_i(f_1, \dots, f_k).
    \quad \forall 1 \leq i \leq k, 1 \leq j \leq d.
\end{align*}
%
On the other hand, when any component $g_i$ is noncommutative,
\cref{eq:polyrec} does not have any number sequence solution.
%
To see this, it suffices to notice that any sequence solution of~\cref{eq:polyrec}
gives rise to a commutative Hadamard-finite series solution of the companion system~\cref{eq:Hadamard-finite companion system}
(via the inverse of the isomorphism above),
but as we have observed~\cref{eq:Hadamard-finite companion system} has a unique, noncommutative solution.

The argument above gives us an algorithm for the consistency problem:
%
Construct the companion system~\cref{eq:Hadamard-finite companion system}
and check that all components of the unique Hadamard-finite series solution are commutative,
which is decidable by~\cref{thm:Hadamard-finite series - decidable commutativity}.
%
\begin{theorem}
    \label{thm:polyrec consistency}
    The polyrec consistency problem~(\cref{problem:polyrec consistency}) is decidable
    with Ackermann complexity.
\end{theorem}
% \begin{proof}
%     Fix the initial condition $c \in \Q^k$.
%     %
%     Thanks to the first part of~\cref{lem:polyrec Hadamard bridge}
%     the companion system has a unique Hadamard-finite series solution with initial condition $c$.
%     %
%     Thanks to the second part of the same lemma,
%     it suffices to check whether all components of this solution are commutative.
%     %
%     The latter can be checked thanks to~\cref{thm:Hadamard-finite series - decidable commutativity}
%     by running $k$ times the decision procedure for commutativity.
% \end{proof}